% Generated draft: requires mathematical and source review.
\documentclass[11pt]{article}
\usepackage[margin=25mm]{geometry}
\usepackage{fontspec}
\setmainfont{Latin Modern Roman}
\usepackage{amsmath,amssymb,mathtools,unicode-math}
\setmathfont{Latin Modern Math}
\usepackage{enumitem,needspace,xurl,hyperref}
\setlength{\parindent}{0pt}
\setlength{\parskip}{5pt}
\setlength{\emergencystretch}{3em}
\newcommand{\N}{\mathbb{N}}
\newcommand{\Z}{\mathbb{Z}}
\newcommand{\Q}{\mathbb{Q}}
\newcommand{\R}{\mathbb{R}}
\newcommand{\C}{\mathbb{C}}
\newcommand{\F}{\mathbb{F}}
\newcommand{\A}{\mathbb{A}}
\newcommand{\PP}{\mathbb{P}}
\newcommand{\E}{\mathbb{E}}
\newcommand{\eps}{\varepsilon}
\newcommand{\dd}{\,\mathrm{d}}
\begin{document}
\section*{UnsolvedMath ID: 20000549}
\textbf{Source problem code:} catalogue entry
\subsection{Definitions and mathematical statement}
A nonzero real polynomial in several variables is real stable if it has no zero when every variable has positive imaginary part; in one variable this is equivalent to all roots being real. A linear stability preserver maps stable polynomials to stable polynomials or zero. The source asks how such operators quantitatively affect root locations, describing this as quantitative versions of MSS. It specifies no operator class, dimension, degree normalization or root statistic, so this is a research direction rather than one universal numerical inequality.

The supplied candidate theorem concerns \(I-tD\), where \(D=d/dx\) and \(t>0\). Let \(p\) be a degree-\(n\) real-rooted polynomial, with largest root \(A\), smallest root \(B\), and diameter \(W=A-B\). The report asserts that \(p-tp'\) has exactly one root above \(A\), which is its largest root. Write it as \(A+u\). For roots \(r_1,\ldots,r_n\) counted with multiplicity, the equivalent displacement equation is
\[
 1=t\sum_{j=1}^n\frac1{u+A-r_j},\qquad u>0.
\]
For \(n\ge2\), its sharp lower and upper extremal displacements can be specified as the positive solutions
\[
 1=t\left(\frac1L+\frac{n-1}{L+W}\right),\qquad
 1=t\left(\frac{n-1}U+\frac1{U+W}\right).
\]
Equivalently these are
\[
 L_{n,t}(W)=\frac{nt-W+\sqrt{(nt-W)^2+4tW}}2,
 \quad
 U_{n,t}(W)=\frac{nt-W+\sqrt{(nt-W)^2+4(n-1)tW}}2.
\]
The reported bound is \(L_{n,t}(W)\le u\le U_{n,t}(W)\). For \(W>0\), its lower extremal root multiset is \(\{A,B,\ldots,B\}\), and its upper one is \(\{A,\ldots,A,B\}\), with precisely these extremizers as supplied. For \(W=0\), all roots coincide and both endpoints are \(nt\); degree one likewise has \(W=0\) and displacement \(t\). The displayed formulas specify the report's endpoints through its stated two-point extremizers; the original summary refers to formulas in unavailable artifacts rather than printing them.

The report also asserts that enlarging the sorted root-gap vector \((A-r_j)_j\) coordinatewise makes the displacement nonincreasing. Separately it uses \(I-aD^2\), \(a>0\), to show that qualitative stability preservation plus fixing constants and linear polynomials cannot imply an operator-independent displacement bound. Thus a fixed-operator sharp theorem is not a bound for every normalized preserver. Finally it mentions an exact Rayleigh-deficit barrier identity linking the univariate calculation to multivariate MSS methods, explicitly as standard synthesis rather than a novelty claim; that identity and its hypotheses are not supplied in the record.
\subsection{Short English statement}
Seek quantitative root-location bounds for stability-preserving operators. The supplied report gives sharp largest-root motion for \(I-tD\), including two-point extremizers and gap monotonicity, and explains why stability plus low-degree normalization alone cannot give a bound uniform over all operators.
\subsection{Sources}
\begin{itemize}
\item Local UnsolvedMath registry, permanent ID \(20000549\), source code \texttt{catalogue entry}; AIM analysis corpus, \texttt{aim-analysis-notes.json}, notes[80]. Attribution: Ulam AI / UnsolvedMath Contributors, CC BY 4.0.
\item AIM \emph{Beyond Kadison--Singer: paving and consequences}, source item b8, attributed in the extracted bullet to ``Petter Branden'': \url{https://docs.google.com/document/d/1XRw0WyG5PshqboCp99tovCRD1GCK0jXZiNLH_nw618c/edit}. The source mentions MSS without a full citation. Original remarks and literature fields are empty.
\item The candidate summary refers to explicit endpoint formulas and a barrier identity in attempt artifacts not reproduced in the supplied record. No authoritative primary citation is supplied for the quantitative theorem.
\item Registry dated August 14, 2026: statement \texttt{reconstructed\_unverified}, status \texttt{partially\_solved}, classification \texttt{PARTIAL-PROGRESS}; attempt 1 labelled partial result, rigor proved, novelty confidence low. The root-motion theorem and obstruction are machine-generated research aids, not peer reviewed, requiring independent verification before citation.
\end{itemize}
\end{document}
